% ============================================================
%  共用宏包与命令 —— 两个 main.tex 都 \input 这个文件
%  作者单位：真实值放在本地 affiliation.tex（已被 .gitignore 拦截）。
%    · 有该文件  → 用真实单位（本机编译）
%    · 无该文件  → 用下面占位值（公开仓库 / 别人 clone 后编译）
% ============================================================

\usepackage[T1]{fontenc}
\usepackage[utf8]{inputenc}
\usepackage{booktabs}       % 三线表
\usepackage{graphicx}
\usepackage{url}
\usepackage{amsmath,amssymb}
\usepackage{booktabs,xcolor,colortbl}
\usepackage{hyperref}
\hypersetup{colorlinks=true,linkcolor=blue,citecolor=blue,urlcolor=blue}

% ---------- 自定义命令（全篇统一，不要在章节里另写） ----------
\newcommand{\agnes}{\textsc{Agnes}}
\newcommand{\comfy}{\textsc{ComfyUI}}
\newcommand{\styledef}{STYLE\_LOCK}
\newcommand{\insight}{InsightFace}
\newcommand{\vds}{VDS}

\newcommand{\nts}{\textbf{S1E01}}

% 行内代码统一风格
\newcommand{\code}[1]{\texttt{#1}}

% 场景代号强调（论文里会反复出现场景号，别写漏下标）
\newcommand{\shot}{\textsc{Shot}}

% ---------- 图表默认参数 ----------
\graphicspath{{figures/}}
\setcounter{topnumber}{3}
\setcounter{bottomnumber}{3}
\setcounter{totalnumber}{6}
\renewcommand{\topfraction}{0.9}
\renewcommand{\bottomfraction}{0.9}

% ---------- 作者 / 单位 ----------
% 真实单位只在本地 affiliation.tex。该文件被 .gitignore 拦截，不会进公开仓库。
% 顺序不能颠倒：先读真值，再用 \providecommand 补占位（已定义的不覆盖）。
\IfFileExists{affiliation.tex}{\input{affiliation.tex}}{}
\providecommand{\theauthor}{[Author name withheld until submission]}
\providecommand{\theaffil}{[Affiliation withheld until submission]}
\providecommand{\theemail}{[email]}

% ---------- 语言与排版 ----------
\hyphenation{ComfyUI InsightFace negative-prompt}
